\documentclass[a4paper,12pt]{article}
\setlength{\headheight}{68.803pt}
\addtolength{\topmargin}{-56.803pt}

\usepackage[utf8]{inputenc}
\usepackage[T1]{fontenc}
\usepackage{lmodern}
\usepackage[nopatch=footnote]{microtype}
\usepackage{fvextra}
\DefineVerbatimEnvironment{Verbatim}{Verbatim}{
  breaklines=true,
  breakanywhere=true,
  frame=single,
  fontsize=\small
}

% margin settings
\usepackage[
  left=1in,
  right=1in,
  top=3cm,
  bottom=1in
]{geometry}

\usepackage[
    colorlinks=true,
    linkcolor=black,
    urlcolor=blue,
    citecolor=blue,
    pdfborder={0 0 0},
    pdfauthor={Massimo Mantineo},
    pdftitle={Server TCP con epoll()},
    pdfsubject={Laboratorio di Reti e Sistemi Distribuiti: HANDS-ON 4},
]{hyperref}

\usepackage{graphicx}
\graphicspath{{../../assets/}{../../assets/ho1/}{../../assets/ho2/}{../../assets/ho3/}{../../assets/ho4/}{../../assets/ho5/}{../../assets/ho6/}{../../assets/ho7/}{../../assets/ho8/}{../../assets/ho10/}}
\usepackage{tikz}
\usetikzlibrary{positioning, arrows.meta}
\usepackage{listings}
\usepackage{xcolor}
\usepackage{amsmath, amssymb}
\usepackage{fancyhdr}
\usepackage{setspace}
\usepackage{pgfplots}
\pgfplotsset{compat=1.17}
\usepackage{booktabs}
\usepackage[skins, listings]{tcolorbox}
\usepackage{float}


\newtcblisting{terminal}[1][]{
    listing only,
    colback=black!85!white,
    colframe=black!70!white,
    colupper=green!80!white,
    coltitle=white,
    title=#1,
    fonttitle=\bfseries\sffamily,
    listing options={
        language={},
        basicstyle=\ttfamily\small\color{green!80!white},
        backgroundcolor={},
        frame=none,
        breaklines=true,
        showstringspaces=false,
        numbers=none
    },
    enhanced,
    drop shadow,
    arc=3mm
}

\newtcblisting{ccode}[1][]{
    listing only,
    colback=blue!3!white,
    colframe=blue!70!black,
    title=#1,
    fonttitle=\bfseries\sffamily,
    listing options={
        style=cstyle,
        backgroundcolor={},
        frame=none
    },
    enhanced,
    drop shadow,
    arc=3mm
}

\setlength{\footskip}{50pt}

% settings for C listings
\lstdefinestyle{cstyle}{
    language=C,
    basicstyle=\ttfamily\small,
    numbers=left,
    numberstyle=\tiny,
    stepnumber=1,
    numbersep=5pt,
    backgroundcolor=\color{white},
    showspaces=false,
    showstringspaces=false,
    showtabs=false,
    frame=single,
    tabsize=4,
    captionpos=b,
    breaklines=true,
    breakatwhitespace=true,
    keywordstyle=\color{blue}\bfseries,
    commentstyle=\color{green!50!black},
    stringstyle=\color{red},
}
\lstset{style=cstyle}

% Header and Footer
\pagestyle{fancy}
\fancyhf{}

\fancyhead[L]{%
  \parbox[b]{0.45\textwidth}{%
    \small
    Student Name: Massimo Mantineo\\
    Student ID: 541924
    \vspace{20pt}
  }%
}
\fancyhead[C]{%
  \parbox[b]{0.2\textwidth}{%
    \centering
    \normalsize \bfseries
    LabReSiD25\\
    Hands-On 4
    \vspace{20pt}
  }%
}
\fancyhead[R]{%
  \parbox[b]{0.35\textwidth}{%
    \flushright
    \includegraphics[width=3.5cm]{\detokenize{logo_unime_orizontale.png}}
    \vspace{15pt}
  }%
}

\renewcommand{\contentsname}{Indice}

\title{\vspace{-2cm}Report: Server TCP ad Alte Prestazioni con \texttt{epoll()}\\
\large (Hands-On 4)}
\author{Massimo Mantineo \\ Matricola: 541924}
\date{MESSINA, A.A. 2024/2025}

\begin{document}

% Cover page
\begin{titlepage}
    \thispagestyle{empty}
    \centering
    \vspace*{1cm}
    \includegraphics[width=6cm]{\detokenize{logo_unime_orizontale.png}}\par\vspace{1cm}
    \vspace{2cm}
    {\Large \textbf{Università degli Studi di Messina}\par}
    {\large Dipartimento MIFT - Corso di Laurea Triennale in Informatica (L-31)\par}
    \vspace*{1cm}
    {\large Laboratorio di Reti e Sistemi Distribuiti\par}
    \vspace{1cm}
    {\Large \textbf{HANDS-ON 4:}\\
    Server TCP ad Alte Prestazioni con \texttt{epoll()}\par}
    \vspace{2cm}
    {\large Massimo Mantineo \par}
    {\large Matricola: 541924 \par}
    \vfill
    {\large Messina, A.A. 2024/2025 \par}
\end{titlepage}

\clearpage
\pagestyle{fancy}
\tableofcontents
\newpage

\section{Problem Statement}
L'ottavo Hands-On (HO4) si focalizza sullo sviluppo di un server TCP ad alte prestazioni su piattaforma Linux. L'obiettivo primario è migrare dal modello di multiplexing tradizionale (\texttt{select}) alla moderna interfaccia \textbf{\texttt{epoll}}. Il sistema deve essere in grado di gestire simultaneamente fino a 10 connessioni client, garantendo l'integrità dei messaggi ricevuti e fornendo un feedback immediato (\texttt{ACK\_EPOLL\_OK}). La sfida tecnica consiste nell'ottimizzare l'uso delle risorse di sistema riducendo l'overhead di scansione dei descrittori tipico delle soluzioni legacy.

\section{Stato dell'Arte}
Il kernel Linux ha introdotto \texttt{epoll} (Event Poll) per superare i limiti di scalabilità della \texttt{select()} e della \texttt{poll()}.
\begin{itemize}
    \item \textbf{Scalabilità O(1)}: A differenza della \texttt{select}, la cui complessità cresce proporzionalmente al numero di FD monitorati ($O(N)$), \texttt{epoll} ha una complessità di risveglio costante $O(k)$, dove $k$ è il numero di eventi attivi.
    \item \textbf{Meccanismo Kernel-Side}: \texttt{epoll} mantiene una struttura dati ad albero (\textit{Red-Black Tree}) nel kernel per tracciare i socket di interesse. Questo evita di copiare l'intero set di descrittori dallo spazio utente a quello kernel a ogni chiamata.
    \item \textbf{Trigger Modes}: 
    \begin{itemize}
        \item \textit{Level Triggered (LT)}: Notifica l'evento finché il buffer non è vuoto.
        \item \textit{Edge Triggered (ET)}: Notifica l'evento solo al momento della transizione (nuovo dato). Richiede l'uso di socket non-bloccanti per evitare stalli.
    \end{itemize}
\end{itemize}

\section{Metodologia ed Implementazione}
L'implementazione sfrutta le tre primitive cardine dell'API epoll per gestire il ciclo di vita delle connessioni.

\begin{figure}[H]
    \centering
    \begin{tikzpicture}[
        node distance=1cm,
        kernel/.style={rectangle, draw=red!70, fill=red!5, thick, minimum width=8cm, minimum height=4cm, dashed, align=center},
        fd/.style={circle, draw=blue!70, fill=blue!10, thick, minimum size=0.8cm, font=\tiny\ttfamily},
        box/.style={rectangle, draw=orange!70, fill=orange!5, thick, minimum width=3cm, minimum height=1cm, font=\small\sffamily\bfseries}
    ]
        % Kernel Space
        \node[kernel] (kspace) {};
        \node[above=0.1cm of kspace.north, red!70, font=\bfseries\sffamily] {KERNEL SPACE (Linux)};
        
        % Red-Black Tree (Implicitly represented by nodes)
        \node[fd] (fd1) at (-2.5,0.5) {fd 3};
        \node[fd] (fd2) at (-1.5,-0.5) {fd 5};
        \node[fd] (fd3) at (-0.5,0.5) {fd 6};
        \node[fd, fill=yellow!30] (fd\_act) at (1.5,0) {fd 8};
        \node[right=0.1cm of fd\_act, font=\tiny, color=red] {EVENT!};
        
        \node[below=1cm of fd2, font=\itshape\small] {Interest List (RB-Tree)};

        % Ready List
        \node[rectangle, draw=green!70, fill=green!5, thick, minimum width=4cm, minimum height=1cm] (ready) at (1.5,-2.5) {Ready List};
        \draw[arrows={-Stealth}, thick, red] (fd\_act) -- (ready);

        % User Space
        \node[box, below=1.5cm of ready] (user) {epoll\_wait()};
        \draw[arrows={Double-Stealth}, thick, green!50!black] (ready) -- node[right, font=\tiny] {O(k) Return} (user);
        
        \node[below=0.1cm of user, font=\bfseries\sffamily] {USER SPACE};
    \end{tikzpicture}
    \caption{Architettura di \texttt{epoll}: Gestione della Interest List ad albero e Ready List lineare.}
    \label{fig:epoll\_arch}
\end{figure}

\subsection{Inizializzazione e Registrazione}
Viene creata un'istanza di epoll e il socket di ascolto viene aggiunto alla lista di interesse:
\begin{ccode}[Snippet Registrazione epoll]
epoll\_fd = epoll\_create1(0);
ev.events = EPOLLIN; // Monitoraggio in lettura
ev.data.fd = listen\_fd;
epoll\_ctl(epoll\_fd, EPOLL\_CTL\_ADD, listen\_fd, &ev);
\end{ccode}

\subsection{Event Loop Asincrono}
Il server rimane in attesa bloccante su \texttt{epoll\_wait}, che restituisce solo i descrittori pronti:
\begin{ccode}[Snippet Loop degli Eventi]
int nfds = epoll\_wait(epoll\_fd, events, MAX\_EVENTS, -1);
for (int n = 0; n < nfds; ++n) {
    if (events[n].data.fd == listen\_fd) {
        // Accettazione e aggiunta nuovo client con EPOLL\_CTL\_ADD
    } else {
        // Lettura dati e invio ACK\_EPOLL\_OK
    }
}
\end{ccode}


\subsection{Istruzioni di Esecuzione}
Compilazione del server basato su epoll:
\begin{terminal}[Build]
$ gcc server\_epoll.c -o server\_epoll
\end{terminal}

\section{Risultati}
La validazione è stata condotta tramite test di carico sintetico e analisi comparativa.

\begin{figure}[H]
    \centering
    \includegraphics[width=0.9\textwidth]{ho4_wireshark\_multi.png}
    \caption{Cattura epoll multi-client su Wireshark.}
    \label{fig:wireshark\_epoll\_multi}
\end{figure}

\subsection{Test di Concorrenza}
L'esecuzione di 5 client simultanei ha confermato la corretta gestione dei File Descriptor. Il server ha riscontrato latenze minime, visualizzando sui log l'interlacciamento degli ID dei socket (fd:5, fd:6, ecc.) senza mostrare segni di stallo.

\subsection{Analisi Wireshark}
L'ispezione della cattura (Figura~\ref{fig:wireshark\_epoll\_multi}) evidenzia che il server \texttt{epoll} risponde istantaneamente a ogni pacchetto \texttt{PSH} ricevuto. La stabilità della connessione è garantita dalla gestione asincrona degli acknowledgment a livello applicativo.

\subsection{Benchmark Prestazionali}
Per valutare l'efficienza di \texttt{epoll} rispetto al modello \texttt{select} (analizzato nell'Hands-On 3), è stato eseguito un test di carico sintetico con 1000 richieste concorrenti.
\begin{terminal}[Selective vs Epoll Benchmark]
$ ./benchmark\_comparison.sh 1000
Testing select (HO3)...
select (HO3) result: 1.42 seconds
Testing epoll (HO4)...
epoll (HO4) result: 0.85 seconds

Conclusion: epoll is ~40% faster under concurrency.
\end{terminal}
I risultati confermano che, all'aumentare del numero di connessioni, la complessità $O(1)$ di \texttt{epoll} garantisce una scalabilità superiore rispetto al modello $O(N)$ di \texttt{select}, riducendo drasticamente il tempo di CPU speso nel monitoraggio dei descrittori inattivi.

\section{Conclusioni e Sviluppi Futuri}
L'Hands-On 4 ha dimostrato la superiorità tecnologica di \texttt{epoll} per la creazione di servizi di rete moderni. La riduzione della copia dei dati tra user-space e kernel rende questa soluzione lo standard per architetture cloud ad alte prestazioni.

\noindent \textbf{Sviluppi futuri:}
\begin{itemize}
    \item \textbf{io\_uring}: Esplorazione del nuovo sottosistema Linux \texttt{io\_uring} per eliminare completamente le syscall bloccanti tramite code di completamento asincrone.
    \item \textbf{Multi-Core}: Implementazione di un modello \textit{multi-threaded epoll} per distribuire la gestione degli eventi su più core CPU.
\end{itemize}

\end{document}
